\documentclass[a4paper,12pt,reqno]{amsart}
\usepackage[T1]{fontenc}
\usepackage[utf8]{inputenc}
\usepackage{amsmath,amssymb,mathtools,microtype}
\usepackage[colorlinks=true,linkcolor=blue,citecolor=blue,urlcolor=blue]{hyperref}
\hypersetup{
 pdftitle={There are no Riesz bases of exponentials in balls and triangles},
 pdfauthor={Joaquim Ortega-Cerdà}
}

\newtheorem{theorem}{Theorem}[section]
\newtheorem{lemma}[theorem]{Lemma}
\newtheorem{proposition}[theorem]{Proposition}
\newtheorem{corollary}[theorem]{Corollary}
\theoremstyle{remark}
\newtheorem{remark}[theorem]{Remark}
\numberwithin{equation}{section}

\newcommand{\R}{\mathbb R}
\newcommand{\C}{\mathbb C}
\newcommand{\Hh}{\mathcal H}
\newcommand{\Ff}{\mathcal F}
\newcommand{\one}{\mathbf 1}
\DeclareMathOperator{\Ran}{Ran}
\DeclareMathOperator{\supp}{supp}

\title[No Riesz bases of exponentials]{There are no Riesz bases of exponentials in balls and triangles}
\author{Joaquim Ortega-Cerdà}
\address{Universitat de Barcelona\\Centre de Recerca Matemàtica (CRM)}
\email{jortega@ub.edu}
% Funding acknowledgment adapted from https://arxiv.org/html/2606.08304v2.
\thanks{The author is supported by grant PID2024-160033NB-I00, funded by
MICIU/AEI/10.13039/501100011033 and FEDER, UE, and by the Generalitat de
Catalunya through grant 2024 ICREA 00142.}
\date{September 13, 2026}
\subjclass[2020]{42C15, 42A55, 46E20}
\keywords{Riesz basis, exponential system, Paley--Wiener space, stationary projections, boundary geometry}

\begin{document}
\raggedbottom

\begin{abstract}
We prove that $L^2(\Omega)$, where $\Omega$ is a ball in $\R^n$, $n\geq2$, or a nondegenerate triangle in $\R^2$, admits no Riesz basis of exponentials.
Using Beurling weak limits of translates of the frequency set, we construct
a stationary system carrying spectral cutoff projections and comparison
projections built from smooth bumps. The Riesz basis assumption forces
the two projections to remain uniformly at operator-norm distance less
than one. The comparison family is continuous in operator norm, while
boundary geometry produces two pointwise limits of the spectral cutoffs
whose ranges are strictly nested. Both limiting projections would then
lie at distance less than one from the same comparison projection, which
is impossible.
\end{abstract}

\maketitle

\section{Introduction and main result}

Fix an integer $n\geq2$.  Let $\Omega\subset\R^n$ be an open ball and, for $\lambda\in\R^n$, write
\[
 e_\lambda(x)=e^{2\pi i\lambda\cdot x}.
\]
We ask whether there can be a sequence $\Lambda\subset\R^n$ for which
$\{e_\lambda\}_{\lambda\in\Lambda}$ is a Riesz basis of $L^2(\Omega)$.  Thus the
synthesis map
\[
 (c_\lambda)\longmapsto \sum_{\lambda\in\Lambda}c_\lambda e_\lambda
\]
is an isomorphism from $\ell^2(\Lambda)$ onto $L^2(\Omega)$; equivalently, there
are constants $0<A\leq B<\infty$ such that
\begin{equation}\label{eq:riesz}
 A\sum_{\lambda\in\Lambda}|c_\lambda|^2
 \leq
 \left\|\sum_{\lambda\in\Lambda}c_\lambda e^{2\pi i\lambda\cdot x}\right\|_{L^2(\Omega)}^2
 \leq
 B\sum_{\lambda\in\Lambda}|c_\lambda|^2
\end{equation}
for every finitely supported family $(c_\lambda)$, and the exponential system
is complete.

The existence of an exponential Riesz basis for Euclidean balls is a longstanding
folklore problem in harmonic analysis.  Bruna, Massaneda and Ortega-Cerdà
\cite[Section~5, p.~350]{BrunaMassanedaOrtegaCerda} explicitly formulated
the conjecture that balls in dimensions greater than one have no such basis.
The planar case was still recorded as open in 2025
\cite[Section~1.2]{LevTselishchev}.

The main result is the following.

\begin{theorem}\label{thm:ball}
For every $n\geq2$, no ball of positive radius in $\R^n$ admits a Riesz basis of exponentials.
\end{theorem}

An invertible affine change of variables gives the same conclusion for
all nondegenerate ellipsoids in $\R^n$; see Section~\ref{sec:ellipsoids}.
The argument also excludes exponential Riesz bases on every nondegenerate
triangle in $\R^2$ (Theorem~\ref{thm:triangle}).  More generally, Remark~\ref{rem:general-boundary} states a
boundary-measure criterion under which the same projection obstruction
applies to other open sets.

The orthogonal version of this problem was initiated by
Fuglede \cite{Fuglede}.  Iosevich, Katz and Pedersen
\cite{IosevichKatzPedersen} proved that balls in dimensions at least two
admit no orthogonal basis of exponentials.  Quantitative restrictions on
orthogonal exponential systems on a planar ball were obtained by Iosevich and
Kolountzakis \cite{IosevichKolountzakis}, and subsequently strengthened by
Zakharov \cite{Zakharov}.  These results concern pairwise orthogonality.
For Riesz bases, Kozma, Nitzan and Olevskii \cite{KozmaNitzanOlevskii}
constructed a bounded subset of the real line without an exponential
Riesz basis and obtained a quantitative lower bound for the Riesz constant
of a hypothetical basis on a planar ball, without excluding its existence.

Positive results are known for polygonal domains.
Lyubarskii and Rashkovskii \cite{LyubarskiiRashkovskii} constructed
exponential Riesz bases for centrally symmetric convex polygons whose
vertices lie on a lattice
using entire functions of two complex variables.
Debernardi and Lev \cite{DebernardiLev} later removed the arithmetic restriction, proving existence
for every centrally symmetric convex polygon and, more generally, for
centrally symmetric convex polytopes whose faces of all dimensions are
centrally symmetric.

A related construction principle is lattice multi-tiling: the lattice
translates of a bounded measurable domain cover almost every point the
same finite number of times.  Grepstad and Lev \cite{GrepstadLev} proved
that such domains admit exponential Riesz bases under a Riemann
measurability assumption; Kolountzakis \cite{Kolountzakis} gave a simpler
proof and removed this regularity assumption.  Kozma and Nitzan proved
existence for finite unions of intervals \cite{KozmaNitzanIntervals}
and for finite unions of rectangles with sides parallel to the coordinate
axes \cite{KozmaNitzanRectangles}.

There is a function-theoretic reformulation of the problem.  If
\[
 PW_\Omega
 =
 \left\{
 F(z)=\int_\Omega f(x)e^{-2\pi iz\cdot x}\,dx : f\in L^2(\Omega)
 \right\},
 \qquad z\in\C^n,
\]
with the norm transported from $L^2(\Omega)$, then $PW_\Omega$ is a Hilbert space of
entire functions.  For real $\lambda$,
\[
 F(\lambda)=\langle f,e_\lambda\rangle_{L^2(\Omega)}.
\]
Hence $\Lambda$ is the frequency set of an exponential Riesz basis precisely
when the restriction map
\[
 F\longmapsto \bigl(F(\lambda)\bigr)_{\lambda\in\Lambda}
\]
is an isomorphism from $PW_\Omega$ onto $\ell^2(\Lambda)$.  This map is the adjoint of synthesis under the
Fourier--Laplace identification.  In other words, the problem is to clarify the existence of complete
interpolating sequences for a Paley--Wiener space in $n$ variables.

The proof relies on a geometric
feature of the ball: a nontrivial translate of its boundary meets the original
boundary in a set of $(n-1)$-dimensional surface measure zero.

The proof has four steps. First, separation of the frequency set allows
us to replace its points by orthonormal translates of a small smooth bump.
The common Fourier factor is invertible on the bounded domain, so the
basis property gives an exact estimate $\|P-Q_\Lambda\|<1$ between two
orthogonal projections on $L^2(\R^n)$. Second, we take weak limits of real
translates of $\Lambda$ and construct an invariant probability measure on
the resulting compact space. This use of weak limits originates in
Beurling's work on sampling and interpolation \cite{Beurling}; a modern
account of the topology is given in \cite[Section~4]{GOR}.
Third, scalar Bochner measures and continuous box averages give spectral
cutoffs $\Pi_t$ and comparison projections $M_t$ on the stationary space,
with $\|\Pi_t-M_t\|\leq\gamma<1$ for almost every $t$.
The averaging estimate is proved by applying the original projection bound
to compactly supported functions and integrating their scalar inner
products. Finally, a boundary crossing gives two limiting
projections with a strict inclusion of their ranges. Norm-continuity of $M_t$ places both at distance at most
$\gamma$ from $M_{t_0}$, a contradiction.

Sections~\ref{sec:bumps}--\ref{sec:distance} establish the analytic
construction for any bounded open domain with boundary of Lebesgue measure
zero. Section~\ref{sec:sphere} supplies the geometric step for balls,
and Sections~\ref{sec:triangles} and \ref{sec:ellipsoids} give the other
applications.

\subsection*{Acknowledgments}
The author thanks Joaquim Bruna for extensive discussions on the manuscript.

\subsection*{Use of artificial intelligence}
A large language model was used for mathematical discussions and assistance
with the Lean formalization.

The Lean formalization is available at
\href{https://github.com/joaquimortega/RieszEuclidean}{RieszEuclidean}.
The repository contains the proof blueprint, the Lean sources, a standalone
proof file, and the verification records. All blueprint obligations are
proved, and the public statements have been checked by Comparator and
Lean's kernel.

\section{From complete interpolation to bump projections}\label{sec:bumps}

Throughout Sections~\ref{sec:bumps}--\ref{sec:distance}, $\Omega$ is a
bounded open subset of $\R^n$ of positive measure. Inner products are
linear in the first variable. We use the unitary Fourier transform
\[
 (\Ff a)(x)=\int_{\R^n}a(\xi)e^{2\pi i\xi\cdot x}\,d\xi,
\]
initially defined on $L^1\cap L^2$, and set
\begin{equation}\label{eq:P}
 P=\Ff^{-1}M_{\one_\Omega}\Ff,\qquad
 p(y)=\int_\Omega e^{-2\pi iy\cdot x}\,dx.
\end{equation}
Thus $P$ is an orthogonal projection on $L^2(\R^n)$ and commutes with
translations. For compactly supported $a\in L^2$, Fourier inversion and
Fubini give $(Pa)(v)=\int p(v-w)a(w)\,dw$ almost everywhere.
Indeed $a\in L^1$, $\Omega$ has finite measure and $p$ is bounded, so this
formula follows directly by interchanging the integrals defining $P$.

\subsection{Distance between orthogonal projections}

For orthogonal projections $P$ and $Q$ on the same Hilbert space,
their operator-norm distance is $\|P-Q\|$.  We use this distance throughout.

\begin{lemma}\label{lem:projection-distance}
Let $P$ and $Q$ be orthogonal projections on a Hilbert space.  The map
\[
 P:\Ran Q\longrightarrow\Ran P
\]
is a bounded isomorphism if and only if
\begin{equation}\label{eq:projection-distance}
 \|P-Q\|<1.
\end{equation}
\end{lemma}

\begin{proof}
Assume first that $\|P-Q\|=\gamma<1$.  If $u\in\Ran Q$, then
\[
 \|(I-P)u\|=\|(Q-P)u\|\leq\gamma\|u\|,
\]
and therefore
\[
 \|Pu\|^2
 =\|u\|^2-\|(I-P)u\|^2
 \geq(1-\gamma^2)\|u\|^2.
\]
This gives injectivity and a lower bound.  To obtain surjectivity, put
$A=P|_{\Ran Q}$ and $B=Q|_{\Ran P}$.  On the closed
subspace $\Ran P$,
\[
 I-AB=P(P-Q).
\]
Thus $\|I-AB\|\leq\gamma<1$, so the Neumann series makes $AB$ invertible.
It follows that $A$ is onto, and the lower bound makes its inverse bounded.

Conversely, suppose $A=P|_{\Ran Q}$ is an isomorphism.  Set
$K=1+\|A^{-1}\|$ and $c=K^{-1}$, so $0<c\leq1$, including when both
ranges are zero.  For $u\in\Ran Q$ we have $\|Au\|\geq c\|u\|$.
For $y\in\Ran P$, take $u=A^{-1}y$.  Orthogonality gives
\[
 \|y\|^2=\langle u,y\rangle=|\langle u,Qy\rangle|
 \leq \|u\|\,\|Qy\|\leq K\|y\|\,\|Qy\|.
\]
For $y\neq0$ we may cancel $\|y\|$; for $y=0$ the conclusion is
immediate.  Hence $\|Qy\|\geq c\|y\|$.  Pythagoras now gives
\[
 \|(I-P)Q\|\leq\sqrt{1-c^2}<1,
 \qquad
 \|(I-Q)P\|\leq\sqrt{1-c^2}<1.
\]
For every $x$,
\[
 (P-Q)x=P(I-Q)x-(I-P)Qx,
\]
and the two terms on the right are orthogonal.  Since $(I-Q)x$ and $Qx$
are also orthogonal,
\[
 \|(P-Q)x\|^2
 \leq
 \max\!\left\{
 \|P(I-Q)\|^2,\|(I-P)Q\|^2
 \right\}\|x\|^2.
\]
Now $\|P(I-Q)\|=\|(I-Q)P\|$, and both numbers are strictly smaller than
one.  This proves \eqref{eq:projection-distance}.
\end{proof}

We shall also use the following immediate consequence.

\begin{lemma}\label{lem:inclusion}
Let $R_-$ and $R_+$ be orthogonal projections with
\[
 \Ran R_-\subsetneq\Ran R_+.
\]
There is no orthogonal projection $M$ such that
\[
 \|R_--M\|<1,
 \qquad
 \|R_+-M\|<1.
\]
\end{lemma}

\begin{proof}
By Lemma~\ref{lem:projection-distance}, $M$ maps each of the two ranges isomorphically onto
$\Ran M$.  Choose $w\in\Ran R_+\setminus\Ran R_-$.  Surjectivity on the
smaller range gives $u\in\Ran R_-$ with $Mu=Mw$.  Then
$0\neq w-u\in\Ran R_+$ and $M(w-u)=0$, contradicting injectivity on
$\Ran R_+$.
\end{proof}

\subsection{Separated frequencies and an exact bump representation}

\begin{lemma}\label{lem:bumps}
Suppose that $\{e_\lambda:\lambda\in\Lambda\}$ is a Riesz basis of
$L^2(\Omega)$. There are $\delta,r>0$, a real nonnegative
$b\in C_c^\infty(\R^n)$ and $\gamma<1$ such that $\Lambda$ is
$\delta$-separated,
\begin{equation}\label{eq:bump}
 \|b\|_2=1,\quad \supp b\subset B(0,r),\quad 2r<\delta,
 \quad \inf_{x\in\overline\Omega}|\Ff b(x)|>0,
\end{equation}
and the orthogonal projection $Q_\Lambda$ onto the closed span of
$\{b(\cdot-\lambda):\lambda\in\Lambda\}$ satisfies
\begin{equation}\label{eq:initial-distance-bound}
 \|P-Q_\Lambda\|\leq\gamma<1.
\end{equation}
\end{lemma}

\begin{proof}
Apply the lower Riesz bound to coefficients $1$ and $-1$ at distinct
frequencies. Since $|e^{iu}-e^{iv}|\leq|u-v|$,
\[
 2A\leq\|e_\lambda-e_\eta\|_{L^2(\Omega)}^2
 \leq4\pi^2|\lambda-\eta|^2\int_\Omega|x|^2\,dx.
\]
The integral is finite and positive. This gives a positive separation
constant $\delta$ and, by disjoint balls of radius $\delta/2$, a uniform
bound on the number of frequencies in any ball of fixed radius.

Choose a real nonnegative smooth compactly supported $b_0$ with
$\|b_0\|_2=1$. For small $\varepsilon>0$ put
$b(\xi)=\varepsilon^{-n/2}b_0(\xi/\varepsilon)$. Its support is as small
as required and
\[
 (\Ff b)(x)=\varepsilon^{n/2}(\Ff b_0)(\varepsilon x).
\]
Since $(\Ff b_0)(0)>0$ and $\overline\Omega$ is compact, this proves
\eqref{eq:bump} for a sufficiently small fixed $\varepsilon$.

The translates $b_\lambda=b(\cdot-\lambda)$ have disjoint supports and
norm one. Consequently
\[
 W_\Lambda:\ell^2(\Lambda)\longrightarrow\Ran Q_\Lambda,
 \qquad W_\Lambda c=\sum_{\lambda\in\Lambda}c_\lambda b_\lambda,
\]
is unitary. For finitely supported $c$,
\begin{equation}\label{eq:bump-synthesis}
 \Ff(PW_\Lambda c)
 =\one_\Omega(\Ff b)\sum_{\lambda\in\Lambda}c_\lambda e_\lambda.
\end{equation}
Both sides extend continuously to $\ell^2(\Lambda)$. On the right,
synthesis is an isomorphism onto $L^2(\Omega)$ and multiplication by
$\Ff b$ is an isomorphism of that space. Hence
$P:\Ran Q_\Lambda\to\Ran P$ is an isomorphism.
Lemma~\ref{lem:projection-distance} gives \eqref{eq:initial-distance-bound}.
\end{proof}

Fix $b,\delta,r$ and $\gamma$ from this lemma. For every $\delta$-separated
set $\Gamma$, let $Q_\Gamma$ be the orthogonal projection onto the closed
span of the corresponding bumps. Explicitly,
\begin{equation}\label{eq:Q}
 Q_\Gamma a=\sum_{\lambda\in\Gamma}\langle a,b_\lambda\rangle b_\lambda.
\end{equation}
The bumps represent the original synthesis map exactly.

\section{Beurling weak limits and a stationary system}\label{sec:weak}

\subsection{Weak convergence of separated sets}

We use Beurling's notion of weak convergence of point sets
\cite{Beurling}. For a modern treatment see \cite[Section~4]{GOR}.
A sequence of closed sets $\Gamma_j$ converges weakly to $\Gamma$ if,
for every $R,\varepsilon>0$, eventually
\begin{equation}\label{eq:weak-sets}
 \Gamma_j\cap B(0,R)\subset\Gamma+B(0,\varepsilon),\qquad
 \Gamma\cap B(0,R)\subset\Gamma_j+B(0,\varepsilon).
\end{equation}
For sets with a common positive separation constant, this is equivalent
to vague convergence of their counting measures
$\eta_\Gamma=\sum_{\lambda\in\Gamma}\delta_\lambda$.
To see this, weak convergence matches the finitely many points on each
compact region, and continuous compactly supported test functions give
vague convergence. Conversely, testing on small disjoint balls around
limit points forces the matching points to exist. A sequence of unmatched
points in a fixed compact set would have a convergent subsequence and
would be detected by another nonnegative continuous test function.
Uniform separation prevents more than one point from approaching a limit
point.

\begin{lemma}\label{lem:hull}
The space of $\delta$-separated subsets of $\R^n$, including the empty
set, is compact and metrizable in the topology \eqref{eq:weak-sets}.
Translations act jointly continuously (in the translation and in the subset) on this space.
\end{lemma}

\begin{proof}
The separation bound gives a uniform bound on the mass of the counting
measures on each compact set. A diagonal extraction against a countable
dense collection of continuous compactly supported functions gives a
vaguely convergent subsequence from any sequence of these measures.
Its limit is again the counting measure of a $\delta$-separated set:
on a ball of diameter less than $\delta$ there is at most one point,
and a point either converges along a subsequence or leaves each compact
subset. Applying this observation on a countable cover identifies the
limit and rules out both fractional masses and multiplicities.
The same countable test functions metrize vague convergence on this
uniformly locally bounded family. Sequential compactness thus implies
compactness. If $\Gamma_j\to\Gamma$ and $y_j\to y$, the local matching
criterion immediately gives $\Gamma_j-y_j\to\Gamma-y$.
\end{proof}

Let
\[
 X=\overline{\{\Lambda-y:y\in\R^n\}},\qquad T_y\Gamma=\Gamma-y.
\]
This is a compact metrizable translation-invariant space. Its elements
are precisely the weak limits of sequences of translates of $\Lambda$.

\subsection{The projection estimate survives weak limits}

We call a sequence of operators $A_j$ pointwise convergent to $A$ if
$\|A_jf-Af\|\to0$ for every fixed vector $f$. We use the same
convention for operator families depending on a real parameter.

\begin{lemma}\label{lem:strong-Q}
If $\Gamma_j\to\Gamma$ in $X$, then $Q_{\Gamma_j}\to Q_\Gamma$
pointwise on $L^2(\R^n)$. In particular,
\begin{equation}\label{eq:hull-distance-bound}
 \|P-Q_\Gamma\|\leq\gamma\qquad(\Gamma\in X).
\end{equation}
\end{lemma}

\begin{proof}
First take a compactly supported $a\in L^2(\R^n)$. Choose a ball
containing the closed $r$-neighborhood of its support in its interior
and having no point of $\Gamma$ on its boundary. Within this ball,
weak convergence and separation give, for large $j$, a bijection
between the finitely many centers of $\Gamma_j$ and $\Gamma$, with
matched centers converging. Only centers in this ball can contribute
to \eqref{eq:Q}. Translations of $b$ are continuous in $L^2$, so each
term of this finite sum converges in $L^2$. Thus
$Q_{\Gamma_j}a\to Q_\Gamma a$. Compactly supported functions are dense,
and all these projections have norm at most one, proving pointwise
convergence.

Conjugating by translations preserves $P$ and sends $Q_\Lambda$ to
the projection for the translated set. Thus every translate satisfies
\eqref{eq:initial-distance-bound}. Taking a pointwise limit and testing
on an arbitrary vector gives \eqref{eq:hull-distance-bound}.
\end{proof}

Since $P$ is nonzero, the empty set cannot belong to $X$: its projection
is zero and would have distance one from $P$. By
Lemma~\ref{lem:projection-distance} and \eqref{eq:bump-synthesis}, each
configuration in $X$ is again the frequency set of an exponential
Riesz basis. Only the uniform projection estimate will be needed below.

\subsection{Translation averaging}

Put $C_R=[-R,R]^n$ and define a probability measure on $X$ by
\[
 \mu_R=\frac1{|C_R|}\int_{C_R}\delta_{T_v\Lambda}\,dv.
\]
Compactness gives a weakly convergent subsequence as $R\to\infty$;
write $\mu$ for its limit (we don't claim it is unique). For $F\in C(X)$ and fixed $y$,
\[
 \left|\int F\circ T_y\,d\mu_R-\int F\,d\mu_R\right|
 \leq\|F\|_\infty\frac{|(C_R+y)\mathbin\triangle C_R|}{|C_R|}
 \longrightarrow0.
\]
Passing to the limit shows that $\mu$ is invariant under every $T_y$.
Let
\begin{equation}\label{eq:U}
 \Hh=L^2(X,\mu),\qquad (U_yf)(\Gamma)=f(\Gamma+y).
\end{equation}
These operators form a unitary representation of $\R^n$. Joint continuity
of the action on the compact space $X$ makes it pointwise continuous on
$C(X)$, and density extends this conclusion to $\Hh$. This Hilbert space
is separable. The constant function $\one$ has norm one and is fixed by
all the $U_y$.

\section{Bochner measures and stationary Fourier projections}\label{sec:bochner}

We construct projections on $\Hh$ that retain the stationary frequencies
$\theta$ for which $t+\theta\in\Omega$. From now on assume in addition
that $|\partial\Omega|=0$. The spectral variables, like the translation
parameters, belong to $\R^n$.

\subsection{Scalar Fourier measures}

For $f\in\Hh$, the continuous function
$y\mapsto\langle U_yf,f\rangle$ is positive definite. Bochner's theorem
\cite[Section~1.4.3, p.~19]{Rudin} gives a unique finite positive
Borel measure $\sigma_f$ on $\R^n$ such that
\begin{equation}\label{eq:bochner}
 \langle U_yf,f\rangle
 =\int_{\R^n}e^{2\pi iy\cdot\theta}\,d\sigma_f(\theta),
 \qquad \sigma_f(\R^n)=\|f\|^2.
\end{equation}
Uniqueness and the parallelogram identity give
\begin{equation}\label{eq:spectral-parallelogram}
 \sigma_{f+g}+\sigma_{f-g}=2\sigma_f+2\sigma_g,
 \qquad \sigma_{cf}=|c|^2\sigma_f.
\end{equation}
Since $U_y\one=\one$, uniqueness also gives
\begin{equation}\label{eq:constant-spectrum}
 \sigma_{\one}=\delta_0.
\end{equation}
The zero-frequency mass will ensure that the boundary crossing changes
the projection nontrivially.

For $a\in L^1(\R^n)$ define
\[
 T_af=\int_{\R^n}a(y)U_yf\,dy,\qquad
 \phi_a(\theta)=\int_{\R^n}a(y)e^{2\pi iy\cdot\theta}\,dy.
\]
The first integral is an $\Hh$-valued integral for each $f$, bounded in
norm by $\|a\|_1\|f\|$. Expanding its squared norm and using
\eqref{eq:bochner} gives
\begin{align*}
 \|T_af\|^2
 &=\int\!\!\int a(y)\overline{a(w)}
       \langle U_{y-w}f,f\rangle\,dy\,dw\\
 &=\int_{\R^n}|\phi_a(\theta)|^2\,d\sigma_f(\theta).
\end{align*}
Fubini applies because $a\in L^1$ and $\sigma_f$ is finite.
The group law and invariance give, for $a,c\in L^1$,
\begin{equation}\label{eq:multiplier-norm}
 \|T_af\|^2=\int|\phi_a|^2\,d\sigma_f,\qquad
 T_aT_c=T_{a*c},\qquad T_a^*=T_{a^*},
\end{equation}
where $a^*(y)=\overline{a(-y)}$. In particular
$\|T_a\|\leq\|\phi_a\|_\infty$. Convolution of kernels corresponds
to multiplication of symbols, and $\phi_{a^*}=\overline{\phi_a}$.
If two kernels have the same symbol, their operators agree by the
norm identity. These facts provide the scalar calculus used below.

\subsection{One measure to control all exceptional sets}

Choose a dense sequence $(h_j)_{j\geq1}$ in the unit sphere of $\Hh$
and set
\begin{equation}\label{eq:control}
 \sigma=\sum_{j=1}^\infty2^{-j}\sigma_{h_j}.
\end{equation}
This is a finite positive measure of mass one. Moreover,
\begin{equation}\label{eq:control-null}
 \sigma(E)=0\quad\Longrightarrow\quad
 \sigma_f(E)=0\quad(f\in\Hh).
\end{equation}
Indeed approximate $f$ by scalar multiples $g$ of the $h_j$.
Then $\sigma_g(E)=0$ and \eqref{eq:spectral-parallelogram} gives
\[
 \sigma_f(E)\leq2\sigma_{f-g}(E)+2\sigma_g(E)
 \leq2\|f-g\|^2\longrightarrow0.
\]
This control measure allows a single choice of boundary point to work
simultaneously for all vectors.

\subsection{Continuous Fej\'er filters}

The normalized overlap of the boxes $C_R=[-R,R]^n$ is
\begin{equation}\label{eq:fejer-weights}
 w_R(y)=\frac{|C_R\cap(C_R-y)|}{|C_R|}
 =\prod_{j=1}^n\left(1-\frac{|y_j|}{2R}\right)_+.
\end{equation}
Put
\[
 K_R(x)=\frac1{|C_R|}\left|\int_{C_R}e^{2\pi iv\cdot x}\,dv\right|^2,
 \qquad m_R=K_R*\one_\Omega.
\]
Plancherel gives $K_R\geq0$ and $\int K_R=1$. The scaling identity
$K_R(x)=R^nK_1(Rx)$ shows that for each $\varepsilon>0$,
$\int_{|x|>\varepsilon}K_R(x)\,dx\to0$. Thus $K_R$ is an approximate
identity. Since $\one_\Omega$ is locally constant away from its boundary,
\begin{equation}\label{eq:fejer-convergence}
 0\leq m_R\leq1,\qquad
 m_R(x)\longrightarrow\one_\Omega(x)\quad(x\notin\partial\Omega).
\end{equation}
Expanding the square in $K_R$ shows that $K_R=\Ff w_R$. Consequently,
with $p$ as in \eqref{eq:P},
\begin{equation}\label{eq:fejer-coefficients}
 m_R(x)=\int_{\R^n}w_R(y)p(y)e^{2\pi iy\cdot x}\,dy.
\end{equation}
This identity follows by Fubini: $w_R$ has compact support and $p$ is
bounded. The same overlap weight will appear in the averaged inner
products in Section~\ref{sec:distance}.

Define
\begin{equation}\label{eq:finite-Pi}
 \Pi_{t,R}f=\int_{\R^n}w_R(y)p(y)e^{2\pi it\cdot y}U_yf\,dy.
\end{equation}
This is $T_a$ for an integrable kernel with symbol $m_R(t+\theta)$.
By \eqref{eq:multiplier-norm} and \eqref{eq:fejer-convergence}, each
$\Pi_{t,R}$ is a self-adjoint contraction.

\subsection{Constructing the cutoff projections}

\begin{proposition}\label{prop:cutoffs}
The set
\begin{equation}\label{eq:good-t}
 \mathcal G=\{t\in\R^n:\sigma(\partial\Omega-t)=0\}
\end{equation}
is measurable and has full Lebesgue measure. For every $t\in\mathcal G$,
$\Pi_{t,R}$ converges pointwise as $R\to\infty$ to an orthogonal
projection $\Pi_t$. For $s,t\in\mathcal G$ and $f\in\Hh$,
\begin{align}
 \|\Pi_tf\|^2
 &=\int_{\R^n}\one_\Omega(t+\theta)\,d\sigma_f(\theta),
 \label{eq:Pi-norm}\\
 \|(\Pi_t-\Pi_s)f\|^2
 &=\int_{\R^n}|\one_\Omega(t+\theta)-\one_\Omega(s+\theta)|^2
       \,d\sigma_f(\theta).
 \label{eq:Pi-difference}
\end{align}
\end{proposition}

\begin{proof}
The function $(t,\theta)\mapsto\one_{\partial\Omega}(t+\theta)$ is
Borel measurable, so its integral against $\sigma$ is measurable.
Tonelli and the boundary-null assumption give
\begin{equation}\label{eq:good-ae}
 \int_{\R^n}\sigma(\partial\Omega-t)\,dt
 =\int_{\R^n}|\partial\Omega-\theta|\,d\sigma(\theta)=0.
\end{equation}
Thus $\mathcal G$ is conull. Fix $t\in\mathcal G$. By
\eqref{eq:control-null}, the boundary exception is null for every
$\sigma_f$. The difference norm formula for the integrable kernels gives
\[
 \|(\Pi_{t,R}-\Pi_{t,S})f\|^2
 =\int|m_R(t+\theta)-m_S(t+\theta)|^2\,d\sigma_f(\theta).
\]
The bounded symbols converge in $L^2(\sigma_f)$ by dominated convergence.
Hence $\Pi_{t,R}f$ is Cauchy and has a limit $\Pi_tf$. Uniform boundedness
makes $\Pi_t$ a bounded linear contraction. Similarly,
\[
 \|(\Pi_{t,R}^2-\Pi_{t,R})f\|^2
 =\int|m_R(t+\theta)^2-m_R(t+\theta)|^2\,d\sigma_f(\theta)
 \longrightarrow0.
\]
Products pass to the pointwise limit for uniformly bounded operators:
\[
 \|\Pi_{t,R}^2f-\Pi_t^2f\|
 \leq\|\Pi_{t,R}\|\,\|\Pi_{t,R}f-\Pi_tf\|
       +\|(\Pi_{t,R}-\Pi_t)\Pi_tf\|\longrightarrow0.
\]
Thus $\Pi_t^2=\Pi_t$. Self-adjointness passes through inner products,
so $\Pi_t$ is orthogonal. Passing to limits in the norm identities for
$\Pi_{t,R}$ and $\Pi_{t,R}-\Pi_{s,R}$ proves the asserted formulas.
\end{proof}

We call $\theta\mapsto\one_\Omega(t+\theta)$ the cutoff symbol of
$\Pi_t$. Here $t$ is a parameter and $\theta$ is a frequency of the
stationary translation action. Formula~\eqref{eq:Pi-norm} says that
$\|\Pi_tf\|^2$ is the spectral mass of $f$ in $\Omega-t$.
Formula~\eqref{eq:Pi-difference} will construct the two boundary limits.

\section{Preserving the distance bound}\label{sec:distance}

We transfer \eqref{eq:hull-distance-bound} to $\Hh$ by averaging scalar
inner products of compactly supported functions on $\R^n$. The comparison
operator is now a family $M_t$. We first prove that its members are
orthogonal projections and that this family is continuous in operator norm.

\subsection{The covariant bump kernel}

The projection $Q_\Gamma$ has the continuous kernel
\begin{equation}\label{eq:q-kernel}
 q_\Gamma(v,w)=\sum_{\lambda\in\Gamma}
 b(v-\lambda)\overline{b(w-\lambda)}.
\end{equation}
For fixed $v$ at most one summand is nonzero. Thus, uniformly in $\Gamma$,
\begin{equation}\label{eq:q-bounds}
 |q_\Gamma(v,w)|\leq\|b\|_\infty^2,\qquad
 q_\Gamma(v,w)=0\quad\text{if }|v-w|>2r.
\end{equation}
Local finiteness, weak convergence and continuity of $b$ imply joint
continuity in $(\Gamma,v,w)$. The kernel satisfies
\begin{align}
 q_{\Gamma+y}(v,w)&=q_\Gamma(v-y,w-y),\label{eq:q-covariance}\\
 q_\Gamma(w,v)&=\overline{q_\Gamma(v,w)},\label{eq:q-hermitian}\\
 \int q_\Gamma(v,u)q_\Gamma(u,w)\,du&=q_\Gamma(v,w).
 \label{eq:q-projection}
\end{align}
The last identity follows by expanding the locally finite sums and using
orthonormality of the bumps. Compact support makes the integrals absolute.

Set $k_y(\Gamma)=q_\Gamma(0,-y)$, and write $B_y$ for multiplication by
$k_y$ on $\Hh$. The map $y\mapsto k_y$ is continuous with values in
$C(X)$: joint continuity is uniform on compact subsets of $X\times\R^n$.
In particular $y\mapsto B_yU_yf$ is continuous for every $f\in\Hh$.
Moreover it vanishes for $|y|>2r$. Define
\begin{equation}\label{eq:M}
 M_tf=\int_{\R^n}e^{2\pi it\cdot y}B_yU_yf\,dy.
\end{equation}
This is an $\Hh$-valued integral with norm at most
$\|f\|\int\|k_y\|_\infty\,dy$.

\begin{proposition}\label{prop:M}
Every $M_t$ is an orthogonal projection. The family is Lipschitz continuous
in operator norm:
\begin{equation}\label{eq:M-continuity}
 \|M_t-M_s\|\leq L|t-s|,\qquad
 L=2\pi\int_{\R^n}|y|\,\|k_y\|_\infty\,dy<\infty.
\end{equation}
\end{proposition}

\begin{proof}
For bounded $f,g$, invariance of $\mu$ and the substitution
$\Gamma'=\Gamma+y$ give
\begin{align*}
 \langle B_yU_yf,g\rangle
 &=\int_X q_\Gamma(0,-y)f(\Gamma+y)\overline{g(\Gamma)}\,d\mu(\Gamma)\\
 &=\int_X q_{\Gamma'}(y,0)f(\Gamma')
                  \overline{g(\Gamma'-y)}\,d\mu(\Gamma')\\
 &=\langle f,B_{-y}U_{-y}g\rangle.
\end{align*}
Density gives $(B_yU_y)^*=B_{-y}U_{-y}$. Taking the adjoint in
\eqref{eq:M} and changing $y$ to $-y$ therefore gives $M_t^*=M_t$.

For idempotence, covariance implies
$k_u(\Gamma+y)=q_\Gamma(-y,-y-u)$. First let $f$ be bounded.
Expanding $M_t^2$ and putting $v=y+u$ yields
\begin{align*}
 (M_t^2f)(\Gamma)
 &=\int e^{2\pi it\cdot v}
       \left[\int q_\Gamma(0,-y)q_\Gamma(-y,-v)\,dy\right]
       f(\Gamma+v)\,dv\\
 &=\int e^{2\pi it\cdot v}q_\Gamma(0,-v)f(\Gamma+v)\,dv
 =(M_tf)(\Gamma),
\end{align*}
where the second equality is \eqref{eq:q-projection}, with the change
of variable $u=-y$. The original double integral has $|y|,|u|\leq2r$
and uniformly bounded coefficients, so Fubini applies. The equality
extends to every $f\in\Hh$ by density and boundedness. Hence $M_t^2=M_t$,
and $M_t$ is an orthogonal projection.

Finally $|e^{2\pi it\cdot y}-e^{2\pi is\cdot y}|\leq2\pi|t-s||y|$.
Applying the triangle inequality to the vector integral proves
\eqref{eq:M-continuity}. Finiteness follows from \eqref{eq:q-bounds}.
\end{proof}

The norm-continuity comes from the integrable, compactly supported kernel.
The translation operators themselves are only required to be pointwise
continuous. This distinction will matter at the boundary crossing.

\subsection{Compactly supported functions from a stationary vector}

Fix $R>0$, $t\in\R^n$, and abbreviate $C=C_R$. For $f\in\Hh$ define,
for almost every $\Gamma$,
\begin{equation}\label{eq:finite-test}
 a_f(\Gamma)(v)=|C|^{-1/2}\one_C(v)e^{-2\pi it\cdot v}f(\Gamma-v).
\end{equation}
These are compactly supported functions of $v\in\R^n$. Tonelli and
stationarity give
\begin{equation}\label{eq:finite-energy}
 \int_X\|a_f(\Gamma)\|_2^2\,d\mu(\Gamma)
 =\frac1{|C|}\int_C\int_X|f(\Gamma-v)|^2\,d\mu(\Gamma)\,dv
 =\|f\|^2.
\end{equation}
In particular $a_f(\Gamma)\in L^2(\R^n)$ for almost every $\Gamma$.
Borel representatives of $f$ give joint measurability, and the same
identity applied to a representative that vanishes almost everywhere
shows that the resulting integrals do not depend on representatives.
Thus the uncountable set of translation parameters causes no null-set
ambiguity.

\subsection{Compute the two averaged inner products}

Let
\begin{equation}\label{eq:finite-M}
 M_{t,R}f=\int w_R(y)e^{2\pi it\cdot y}B_yU_yf\,dy.
\end{equation}
For bounded $f,g$, expansion of the compact spatial integrals gives
\begin{align*}
 &\int_X\langle Q_\Gamma a_f(\Gamma),a_g(\Gamma)\rangle\,d\mu(\Gamma)\\
 &\quad=\frac1{|C|}\int_{C\times C}e^{2\pi it\cdot(v-w)}
   \int_X q_\Gamma(v,w)f(\Gamma-w)
                   \overline{g(\Gamma-v)}\,d\mu(\Gamma)\,dv\,dw.
\end{align*}
In the inner integral substitute $\Gamma'=\Gamma-v$. For $y=v-w$,
covariance turns the kernel into $q_{\Gamma'}(0,-y)$, and the two
functions become $f(\Gamma'+y)$ and $g(\Gamma')$. The inner integral
is therefore $\langle B_yU_yf,g\rangle$. For every integrable function
$H$ of the difference variable,
\[
 \frac1{|C|}\int_{C\times C}H(v-w)\,dv\,dw
 =\int w_R(y)H(y)\,dy.
\]
Consequently
\begin{equation}\label{eq:average-Q}
 \int_X\langle Q_\Gamma a_f,a_g\rangle\,d\mu
 =\langle M_{t,R}f,g\rangle.
\end{equation}

For $P$ its kernel $p(v-w)$ is independent of $\Gamma$. The same
calculation, now with inner integral $\langle U_{v-w}f,g\rangle$, gives
\begin{align}
 \int_X\langle Pa_f,a_g\rangle\,d\mu
 &=\int w_R(y)p(y)e^{2\pi it\cdot y}\langle U_yf,g\rangle\,dy\notag\\
 &=\langle\Pi_{t,R}f,g\rangle.
 \label{eq:average-P}
\end{align}
All kernels are bounded on $C\times C$. Thus the expansions for bounded
$f,g$ are justified by absolute integration. For general $f,g\in\Hh$,
Cauchy--Schwarz, contractivity of $P,Q_\Gamma$ and
\eqref{eq:finite-energy} bound each averaged form by $\|f\|\|g\|$.
The identities extend by density.

\subsection{Apply the distance bound and pass to the limit}

\begin{proposition}\label{prop:distance}
For every $t\in\mathcal G$,
\begin{equation}\label{eq:stationary-distance-bound}
 \|\Pi_t-M_t\|\leq\gamma<1.
\end{equation}
\end{proposition}

\begin{proof}
Subtracting \eqref{eq:average-Q} from \eqref{eq:average-P} gives
\begin{equation}\label{eq:compression-identity}
 \langle(\Pi_{t,R}-M_{t,R})f,g\rangle
 =\int_X\langle(P-Q_\Gamma)a_f(\Gamma),a_g(\Gamma)\rangle\,d\mu.
\end{equation}
Apply \eqref{eq:hull-distance-bound} for almost every $\Gamma$, followed
by Cauchy--Schwarz in $X$ and \eqref{eq:finite-energy}:
\begin{align*}
 |\langle(\Pi_{t,R}-M_{t,R})f,g\rangle|
 &\leq\gamma\int_X\|a_f(\Gamma)\|_2\|a_g(\Gamma)\|_2\,d\mu\\
 &\leq\gamma\|f\|\|g\|.
\end{align*}
Thus $\|\Pi_{t,R}-M_{t,R}\|\leq\gamma$ for every $t$ and $R$.
By \eqref{eq:q-bounds} and dominated convergence,
\begin{equation}\label{eq:M-norm-limit}
 \|M_{t,R}-M_t\|
 \leq\int_{|y|\leq2r}|1-w_R(y)|\,\|k_y\|_\infty\,dy
 \longrightarrow0,
\end{equation}
uniformly in $t$. For $t\in\mathcal G$, Proposition~\ref{prop:cutoffs}
gives pointwise convergence $\Pi_{t,R}\to\Pi_t$. Testing the distance
bound on each vector and taking the limit proves
\eqref{eq:stationary-distance-bound}.
\end{proof}

The finite averages $M_{t,R}$ need not be projections. Their norm limit
$M_t$ is a projection by Proposition~\ref{prop:M}. At a boundary parameter,
the spectral cutoffs will have two different pointwise limits, whereas
$M_t$ has a single operator-norm limit. This gives a common comparison
projection for the two sides.

\section{Crossing the sphere}\label{sec:sphere}

We now prove Theorem~\ref{thm:ball}. Suppose that a ball
$\Omega=B(a,\rho)\subset\R^n$, $n\geq2$, admits an exponential Riesz
basis. The preceding sections give projections $\Pi_t$ and $M_t$ with
\eqref{eq:stationary-distance-bound} for $t\in\mathcal G$. We will find
two boundary limits whose ranges satisfy a strict inclusion, both within distance
$\gamma<1$ of $M_{t_0}$.

Recall that $\Pi_t$ retains the stationary frequencies $\theta$ for
which $t+\theta\in\Omega$. We approach a boundary point $t_0$ from
outside and inside. At $\theta=0$, this changes the cutoff from zero
to one. The geometric step chooses $t_0$ so that every other frequency
on the translated boundary lies in a common spectral null set.

\subsection{Choose a boundary point with no other spectral boundary mass}

Write $S=\partial\Omega$ and $\nu=\mathcal H^{n-1}|_S$. This is a finite
nonzero measure.

\begin{lemma}\label{lem:sphere}
For every nonzero $\theta\in\R^n$,
\begin{equation}\label{eq:sphere-intersection}
 \nu\bigl(S\cap(S-\theta)\bigr)=0.
\end{equation}
Moreover, $S$ has Lebesgue measure zero.
\end{lemma}

\begin{proof}
If $x\in S\cap(S-\theta)$, then $|x-a|^2=\rho^2$ and
$|x+\theta-a|^2=\rho^2$. Subtraction gives
\[
 2(x-a)\cdot\theta=-|\theta|^2.
\]
For $\theta\ne0$, this is a hyperplane. Its intersection with the sphere
is empty, a point, or an $(n-2)$-dimensional sphere, hence has zero
$\mathcal H^{n-1}$-measure. Finally a sphere in $\R^n$ has zero
$n$-dimensional Lebesgue measure.
\end{proof}

We need the reverse viewpoint: fix a boundary point $t_0$ and control
the frequencies $\theta$ for which $t_0+\theta\in S$.
Fubini interchanges these roles.  Using the control measure $\sigma$
from \eqref{eq:control}, we obtain
\begin{equation}\label{eq:boundary-fubini}
 \begin{aligned}
 \int_{S}\sigma\bigl((S-t)\setminus\{0\}\bigr)\,d\nu(t)
 &=\int_{\R^n\setminus\{0\}}
     \nu\bigl(S\cap(S-\theta)\bigr)\,d\sigma(\theta)\\
 &=0.
 \end{aligned}
\end{equation}
The integrand on the left is nonnegative and $\nu(S)>0$.  Thus we can
choose $t_0\in S$ such that
\begin{equation}\label{eq:boundary-null}
 \sigma\bigl((S-t_0)\setminus\{0\}\bigr)=0.
\end{equation}
By \eqref{eq:control-null}, the same set is null for every $\sigma_f$.
This is why we introduced one control measure: the chosen point works
for every vector in $\Hh$ simultaneously.

Zero must be excluded from this null-set statement.  Indeed,
$\sigma_{\one}=\delta_0$, so the constant vector has spectral mass
there.  This mass is precisely what will enter the cutoff.

\subsection{Approach the point from outside and inside}

The good set $\mathcal G$ has full Lebesgue measure in $\R^n$, but this
does not imply that it contains $t_0$.  In fact $t_0\notin\mathcal G$:
since $0\in S-t_0$ and $\sigma_{\one}=\delta_0$, condition
\eqref{eq:good-t} cannot hold there.  We therefore approach $t_0$
through good parameters and take limits of their projections.

Every neighborhood of a point on the sphere meets both $\Omega$ and
its strict exterior in nonempty open sets.  Each has positive Lebesgue
measure, so each meets $\mathcal G$.  We can consequently choose
\[
 t_j^-\in\mathcal G\cap(\R^n\setminus\overline\Omega),\qquad
 t_j^+\in\mathcal G\cap\Omega,\qquad t_j^\pm\longrightarrow t_0.
\]
These choices use open neighborhoods in the ambient space.  There is
no need to restrict the good parameters to a prescribed radial line.

Write the corresponding cutoff symbols as
\[
 a_j^\pm(\theta)=\one_\Omega(t_j^\pm+\theta).
\]
If $t_0+\theta$ is away from $S$, membership in $\Omega$ is locally
constant, so both symbols eventually equal
$\one_\Omega(t_0+\theta)$.  At $\theta=0$, the exterior symbol is
always zero and the interior symbol is always one.  All remaining
cases lie in $(S-t_0)\setminus\{0\}$, which is null for every
$\sigma_f$ by \eqref{eq:boundary-null}.
Thus, for every $f\in\Hh$, we have $a_j^\pm\to m_\pm$
$\sigma_f$-almost everywhere, where
\begin{align}
 m_-(\theta)&=\one_{\Omega-t_0}(\theta),\label{eq:mminus}\\
 m_+(\theta)&=\one_{\Omega-t_0}(\theta)+\one_{\{0\}}(\theta).
 \label{eq:mplus}
\end{align}
Since $\Omega$ is open and $t_0\in S$, zero does not belong to
$\Omega-t_0$.  Both limits are therefore indicator functions, and
$m_-\leq m_+$.  The only added spectral component is at zero.

\subsection{Turn the limiting symbols into projections}

We now justify that the projections themselves converge.  For either
choice of sign, \eqref{eq:Pi-difference} gives
\[
 \|(\Pi_{t_j^\pm}-\Pi_{t_k^\pm})f\|^2
 =\int_{\R^n}|a_j^\pm-a_k^\pm|^2\,d\sigma_f.
\]
The right-hand side tends to zero, since
\[
 |a_j^\pm-a_k^\pm|^2
 \leq2|a_j^\pm-m_\pm|^2+2|a_k^\pm-m_\pm|^2
\]
and dominated convergence applies to both terms.  Hence
$(\Pi_{t_j^\pm}f)_j$ is Cauchy for each $f$.  The uniform bound
$\|\Pi_t\|\leq1$ gives bounded linear operators $R_-,R_+$ with
\begin{equation}\label{eq:boundary-strong}
 \Pi_{t_j^-}\longrightarrow R_-,\qquad
 \Pi_{t_j^+}\longrightarrow R_+
 \quad\text{pointwise}.
\end{equation}
They are orthogonal projections.  To see the idempotence explicitly,
if $P_j=\Pi_{t_j^\pm}$ and $R=R_\pm$, then
\[
 \|P_j^2f-R^2f\|
 \leq\|P_j\|\,\|P_jf-Rf\|+\|(P_j-R)Rf\|\longrightarrow0.
\]
Since $P_j^2=P_j$, we obtain $R^2=R$.  Self-adjointness passes to the
limit through inner products.  Finally, passing to the limit in
\eqref{eq:Pi-norm} gives
\begin{equation}\label{eq:boundary-norm}
 \|R_\pm f\|^2=\int_{\R^n}m_\pm(\theta)\,d\sigma_f(\theta).
\end{equation}
This formula lets us read the order of the limiting projections from
the order of their symbols.

\subsection{Strict inclusion of subspaces and the contradiction}

Since $m_-\leq m_+$, \eqref{eq:boundary-norm} implies
$\|R_-f\|\leq\|R_+f\|$ for every $f$.  In particular,
$R_+f=0$ implies $R_-f=0$.  Thus
\[
 \ker R_+\subset\ker R_-,\qquad
 \Ran R_-\subset\Ran R_+,
\]
where the second inclusion follows by taking orthogonal complements.

To see that this inclusion is strict, use the constant vector.  Its
spectral measure is $\delta_0$, so \eqref{eq:boundary-norm} gives
\[
 R_-\one=0,\qquad \|R_+\one\|=\|\one\|=1.
\]
An orthogonal projection preserves a vector's norm only when it fixes
that vector, by Pythagoras.  Hence $R_+\one=\one$: the constant
function belongs to the larger range and is orthogonal to the smaller
one.  We have proved
\begin{equation}\label{eq:strict-inclusion}
 \Ran R_-\subsetneq\Ran R_+.
\end{equation}
More precisely, the difference is an orthogonal projection, and
\begin{equation}\label{eq:zero-frequency-jump}
 \|(R_+-R_-)f\|^2=\sigma_f(\{0\}),\qquad
 (R_+-R_-)\one=\one.
\end{equation}
The jump need not have rank one; the constant vector ensures that it
is nonzero.


Both limits also inherit a common distance bound. By
Proposition~\ref{prop:M}, $M_{t_j^\pm}\to M_{t_0}$ in operator norm.
For every $f\in\Hh$,
\begin{align*}
 \|(R_\pm-M_{t_0})f\|
 &=\lim_j\|(\Pi_{t_j^\pm}-M_{t_0})f\|\\
 &\leq\limsup_j\left(\gamma+\|M_{t_j^\pm}-M_{t_0}\|\right)\|f\|
 =\gamma\|f\|.
\end{align*}
Consequently
\begin{equation}\label{eq:boundary-distance}
 \|R_--M_{t_0}\|\leq\gamma<1,\qquad
 \|R_+-M_{t_0}\|\leq\gamma<1.
\end{equation}
The comparison at $t_0$ is legitimate even though $t_0\notin\mathcal G$:
$M_t$ is defined and is a projection for every real parameter.
Lemma~\ref{lem:inclusion} contradicts \eqref{eq:strict-inclusion} and
\eqref{eq:boundary-distance}. This proves Theorem~\ref{thm:ball}.

\begin{remark}[The role of the stationary space]\label{rem:stationary-space}
The invariant probability measure supplies a nonzero constant vector
with spectral measure $\delta_0$. The ordinary translation representation
on $L^2(\R^n)$ has no such vector. Its scalar spectral measures are
$|\Ff a(\theta)|^2\,d\theta$, so changing a cutoff at the single point
zero has no effect on that space. On $\Hh$, the same change gives
$R_-\one=0$ and $R_+\one=\one$. The averaging argument brings both this
invariant vector and the original distance bound into the same stationary
space.
\end{remark}

\begin{remark}\label{rem:interval}
For a sphere the boundary has no positive-measure overlap with a
nontrivial translate. This fails for the one-dimensional interval:
a translation can send one endpoint to the other. It also fails for
opposite sides of a rectangle. A boundary crossing can then exchange
spectral components, and the contradiction from strict inclusion does not follow.
\end{remark}

\begin{remark}[A more general boundary condition]\label{rem:general-boundary}
Let $\Omega\subset\R^n$ be a nonempty bounded open set with
$|\partial\Omega|=0$. Suppose its boundary carries a finite nonzero
positive Borel measure $\nu$ such that
\[
 \nu\bigl(\partial\Omega\cap(\partial\Omega-\theta)\bigr)=0
 \qquad(\theta\ne0).
\]
Assume also that, for $\nu$-almost every $t\in\partial\Omega$, every
neighborhood of $t$ meets both $\Omega$ and
$\R^n\setminus\overline\Omega$ in sets of positive Lebesgue measure.
Then $L^2(\Omega)$ admits no Riesz basis of exponentials.
Indeed Sections~\ref{sec:bumps}--\ref{sec:distance} apply under these
hypotheses. Tonelli with $\nu$ supplies a point satisfying both the
spectral boundary-null condition and the local separation condition.
The latter supplies the two sequences of good parameters. Their symbols
have exactly the limits \eqref{eq:mminus}--\eqref{eq:mplus}, and the same
arguments using pointwise limits and norm continuity give a contradiction.
In particular one may use surface measure whenever it is finite and
nonzero and satisfies the displayed overlap and separation conditions.
\end{remark}

\section{Triangles}\label{sec:triangles}

For a triangle, translations along an edge can overlap it in positive
length. The other edges are not parallel to it, so all surviving boundary
components can still be made to enter the cutoff together.

\begin{theorem}\label{thm:triangle}
No nondegenerate triangle in $\R^2$ admits a Riesz basis of exponentials.
\end{theorem}

\begin{proof}
An invertible affine change of variables preserves the Riesz basis
property: frequencies are transformed by the transpose of the linear
part, and the change introduces only unit phases and a fixed norm factor.
Thus we may assume
\[
 \Omega=\{x_1>0,\ x_2>0,\ x_1+x_2<1\}.
\]
Suppose it has an exponential Riesz basis. Since $\Omega$ is bounded
and its boundary has area zero, Sections~\ref{sec:bumps}--\ref{sec:distance}
give the control measure $\sigma$, the conull set $\mathcal G$,
and the projections with \eqref{eq:stationary-distance-bound}.

Let $L=\{(a,0):0<a<1\}$ and $F=\partial\Omega\setminus L$.
The set $F$ consists of the two other closed edges, including all
vertices. For $\nu=\mathcal H^1|_L$ and every $\theta\in\R^2$,
\[
 \nu\bigl(L\cap(F-\theta)\bigr)=0:
\]
a horizontal segment meets a translated vertical or diagonal supporting
line in at most one point. Tonelli therefore gives
\[
 \int_L\sigma(F-t)\,d\nu(t)
 =\int_{\R^2}\nu\bigl(L\cap(F-\theta)\bigr)\,d\sigma(\theta)=0.
\]
Choose $t_0\in L$ with $\sigma(F-t_0)=0$. This set is null for every
$\sigma_f$ by \eqref{eq:control-null}.

Near $t_0$, the triangle coincides with the half-plane $x_2>0$.
Choose good parameters $t_j^+\in\mathcal G\cap\Omega$ and
$t_j^-\in\mathcal G\cap(\R^2\setminus\overline\Omega)$ tending to
$t_0$ inside such a neighborhood. Their second coordinates are positive
and negative, respectively. For each point $t_0+\theta\in L$, membership
in the triangle is likewise locally exactly positivity of the second
coordinate. Therefore
$\one_\Omega(t_j^++\theta)\to1$ and
$\one_\Omega(t_j^-+\theta)\to0$ on $L-t_0$.
Away from the boundary both symbols converge to
$\one_\Omega(t_0+\theta)$. All other boundary points belong to the
common null set $F-t_0$.

The limiting symbols are consequently
\[
 m_-=\one_{\Omega-t_0},\qquad
 m_+=\one_{\Omega-t_0}+\one_{L-t_0}.
\]
They are indicators of sets ordered by inclusion. Formula~\eqref{eq:Pi-difference}
and dominated convergence make $(\Pi_{t_j^\pm}f)_j$ Cauchy for every
$f$. As in Section~\ref{sec:sphere}, the resulting pointwise limits
$R_-,R_+$ are orthogonal projections with
\[
 \|R_\pm f\|^2=\int m_\pm\,d\sigma_f.
\]
Thus $\|R_-f\|\leq\|R_+f\|$ for every $f$, which gives
$\Ran R_-\subset\Ran R_+$ by taking kernels and orthogonal complements.
Since $0\in L-t_0$ and $\sigma_{\one}=\delta_0$, we have
$R_-\one=0$ and $R_+\one=\one$; the inclusion is strict.
Finally $M_{t_j^\pm}\to M_{t_0}$ in norm, so
\eqref{eq:stationary-distance-bound} gives
$\|R_\pm-M_{t_0}\|\leq\gamma<1$, contradicting
Lemma~\ref{lem:inclusion}.

Here $\|(R_+-R_-)f\|^2=\sigma_f(L-t_0)$. The jump may contain nonzero
tangential frequencies; the constant vector guarantees that it is nonzero.
\end{proof}

\begin{remark}[Convex polygons with an unpaired edge]\label{rem:odd-polygons}
The same proof excludes an exponential Riesz basis on the interior of a
convex polygon having an edge with no parallel partner. Edges are taken
to be maximal boundary segments. In particular the conclusion holds for
every convex polygon with an odd number of sides: a bounded convex polygon
has at most two maximal edges in each unoriented direction, so an odd
number of sides cannot all be paired by parallelism.

To verify the extension, let $L$ be the relative interior of an unpaired
edge and put $F=\partial\Omega\setminus L$. Every other edge is
nonparallel to $L$, so for $\nu=\mathcal H^1|_L$,
\[
 \nu\bigl(L\cap(F-\theta)\bigr)=0\qquad(\theta\in\R^2).
\]
The intersections with the translated supporting lines are finite;
the two endpoints contribute only finitely many further points.
Tonelli gives $t_0\in L$ with $\sigma(F-t_0)=0$.
Fix the inward unit normal to $L$. At each point of $L$, membership in
$\Omega$ is locally determined by the positive sign of the displacement
in this normal direction. Choose good parameters approaching $t_0$
from the two local half-planes. For every $t_0+\theta\in L$, the same
normal sign determines their eventual membership. The symbols are
therefore the same $m_-,m_+$ as in the triangle proof, with this edge $L$.
Their limiting projections satisfy $\Ran R_-\subsetneq\Ran R_+$
because $0\in L-t_0$.
Norm-continuity of $M_t$ again supplies the common comparison projection
$M_{t_0}$ and Lemma~\ref{lem:inclusion} gives the contradiction.
\end{remark}

\section{Ellipsoids}\label{sec:ellipsoids}

\begin{corollary}\label{cor:ellipsoid}
For every $n\geq2$, no nondegenerate ellipsoid in $\R^n$ admits a
Riesz basis of exponentials.
\end{corollary}

\begin{proof}
Write the ellipsoid as $\Omega=a+A B(0,1)$, where $A$ is an invertible real
$n\times n$ matrix.  The pullback $f\mapsto[y\mapsto f(a+Ay)]$
is an isomorphism from $L^2(\Omega)$ onto $L^2(B(0,1))$ with norm factor
$|\det A|^{-1/2}$.  It sends $e_\lambda$ to
$e^{2\pi i\lambda\cdot a}e_{A^{\mathsf T}\lambda}$.  The invertible change of
frequency and the unit phases preserve the Riesz basis property.
A basis on $\Omega$ would therefore give one on the unit ball, contradicting
Theorem~\ref{thm:ball}.
\end{proof}

\begin{thebibliography}{99}

\bibitem{Beurling}
A.~Beurling,
\emph{The Collected Works of Arne Beurling}, Vol.~2: Harmonic Analysis,
edited by L.~Carleson, P.~Malliavin, J.~Neuberger and J.~Wermer,
Contemporary Mathematicians, Birkh\"auser, Boston, 1989.
See the Mittag--Leffler lectures on balayage and interpolation.

\bibitem{BrunaMassanedaOrtegaCerda}
J.~Bruna, X.~Massaneda and J.~Ortega-Cerdà,
\emph{Connections between signal processing and complex analysis},
Contrib. Sci. \textbf{2} (2003), no.~3, 345--357.
\href{https://publicacions.iec.cat/repository/pdf/00000022/00000096.pdf}{Publisher's full text}.

\bibitem{DebernardiLev}
A.~Debernardi and N.~Lev,
\emph{Riesz bases of exponentials for convex polytopes with symmetric faces},
J.~Eur. Math. Soc. \textbf{24} (2022), no.~8, 3017--3029.
\href{https://doi.org/10.4171/JEMS/1158}{doi:10.4171/JEMS/1158}.

\bibitem{Fuglede}
B.~Fuglede,
\emph{Commuting self-adjoint partial differential operators and a group
theoretic problem},
J.~Funct. Anal. \textbf{16} (1974), no.~1, 101--121.
\href{https://doi.org/10.1016/0022-1236(74)90072-X}{doi:10.1016/0022-1236(74)90072-X}.

\bibitem{GrepstadLev}
S.~Grepstad and N.~Lev,
\emph{Multi-tiling and Riesz bases},
Adv. Math. \textbf{252} (2014), 1--6.
\href{https://doi.org/10.1016/j.aim.2013.10.019}{doi:10.1016/j.aim.2013.10.019}.

\bibitem{GOR}
K.~Gr\"ochenig, J.~Ortega-Cerd\`a and J.~L.~Romero,
\emph{Deformation of Gabor systems},
Adv. Math. \textbf{277} (2015), 388--425.
\href{https://doi.org/10.1016/j.aim.2015.01.019}{doi:10.1016/j.aim.2015.01.019}.
\href{https://arxiv.org/abs/1311.3861}{arXiv:1311.3861}.

\bibitem{IosevichKatzPedersen}
A.~Iosevich, N.~Katz and S.~Pedersen,
\emph{Fourier bases and a distance problem of Erd\H{o}s},
Math. Res. Lett. \textbf{6} (1999), no.~2, 251--255.
\href{https://doi.org/10.4310/MRL.1999.v6.n2.a13}{doi:10.4310/\allowbreak MRL.1999.v6.n2.a13}.

\bibitem{IosevichKolountzakis}
A.~Iosevich and M.~N.~Kolountzakis,
\emph{Size of orthogonal sets of exponentials for the disk},
Rev. Mat. Iberoam. \textbf{29} (2013), no.~2, 739--747.
\href{https://doi.org/10.4171/RMI/737}{doi:10.4171/RMI/737}.

\bibitem{Kolountzakis}
M.~N.~Kolountzakis,
\emph{Multiple lattice tiles and Riesz bases of exponentials},
Proc. Amer. Math. Soc. \textbf{143} (2015), no.~2, 741--747.
\href{https://doi.org/10.1090/S0002-9939-2014-12310-0}{doi:10.1090/S0002-9939-2014-12310-0}.

\bibitem{KozmaNitzanIntervals}
G.~Kozma and S.~Nitzan,
\emph{Combining Riesz bases},
Invent. Math. \textbf{199} (2015), no.~1, 267--285.
\href{https://doi.org/10.1007/s00222-014-0522-3}{doi:10.1007/s00222-014-0522-3}.

\bibitem{KozmaNitzanRectangles}
G.~Kozma and S.~Nitzan,
\emph{Combining Riesz bases in $\R^d$},
Rev. Mat. Iberoam. \textbf{32} (2016), no.~4, 1393--1406.
\href{https://doi.org/10.4171/RMI/922}{doi:10.4171/RMI/922}.

\bibitem{KozmaNitzanOlevskii}
G.~Kozma, S.~Nitzan and A.~Olevskii,
\emph{A set with no Riesz basis of exponentials},
Rev. Mat. Iberoam. \textbf{39} (2023), no.~6, 2007--2016.
\href{https://doi.org/10.4171/RMI/1411}{doi:10.4171/RMI/1411}.

\bibitem{LevTselishchev}
N.~Lev and A.~Tselishchev,
\emph{Unconditional Schauder frames of exponentials and of uniformly bounded
functions in $L^p$ spaces},
Int. Math. Res. Not. IMRN \textbf{2025} (2025), no.~19, rnaf299.
\href{https://doi.org/10.1093/imrn/rnaf299}{doi:10.1093/imrn/rnaf299}.

\bibitem{LyubarskiiRashkovskii}
Y.~I.~Lyubarskii and A.~Rashkovskii,
\emph{Complete interpolating sequences for Fourier transforms supported by
convex symmetric polygons},
Ark. Mat. \textbf{38} (2000), no.~1, 139--170.
\href{https://doi.org/10.1007/BF02384495}{doi:10.1007/BF02384495}.

\bibitem{Rudin}
W.~Rudin,
\emph{Fourier Analysis on Groups},
Interscience Tracts in Pure and Applied Mathematics, vol.~12,
Interscience Publishers, New York, 1962.
See Section~1.4.3, p.~19 (Bochner's theorem).
\href{https://doi.org/10.1002/9781118165621}{Publisher's edition}.

\bibitem{Zakharov}
D.~Zakharov,
\emph{On sets of orthogonal exponentials on the disk},
preprint, 2024.
\href{https://arxiv.org/abs/2405.14063}{arXiv:2405.14063}.
\end{thebibliography}

\end{document}
